%%%%%%%%%%%%%%%%%%%%%%%%%%%%%%%%%%%%%%%%%%%%%%%%%%%%%%%%%%%%%%%%%%%%%%
%                                                                    %
% Two-weight norm inequalities for the Ces\`aro means of             %
% generalized Hermite expansions,                                    %
% by                                                                 %
% \'Oscar Ciaurri and Juan L. Varona                                 %
%                                                                    %
% DATE: first version, 20-March-2003                                 %
% DATE: last modification, 24-March-2004                             %
%                                                                    %
% Written with AMS-LaTeX 2.0                                         %
%                                                                    %
%%%%%%%%%%%%%%%%%%%%%%%%%%%%%%%%%%%%%%%%%%%%%%%%%%%%%%%%%%%%%%%%%%%%%%

\documentclass{amsart}

%%%%%%%% A LOT OF BIG SEVERAL-LINE EQUATIONS (aligns, etc.).
%%%%%%%% THUS, WE ALLOW TO BREAK THEM IN PAGES.
\allowdisplaybreaks

%\date{March 24, 2003}

\date{PUBLISHED IN \textit{J.\ Comput.\ Appl.\ Math.}\ \textbf{178}
  (2005), no.~1-2, 99--110.}

\theoremstyle{plain}
\newtheorem{Theo}{Theorem}
\newtheorem{Prop}{Proposition}
\newtheorem{Lem}{Lemma}

\theoremstyle{definition}
\newtheorem{Def}{Definition}

\theoremstyle{remark}
\newtheorem{Rem}{Remark}

%%%%%%%%

% Cesaro-Hermite kernel and means:
\newcommand{\cheker}[5]{\mathcal{K}_{#1,#2}^{\{#3\}}%
  (#4,#5)} % Cesaro-Hermite kernel
\newcommand{\chem}[3]{\sigma_{#1,#2}^{\{#3\}}} % Cesaro-Hermite mean
\newcommand{\chemmod}[3]{\widetilde{\sigma}_{#1,#2}^{\{#3\}}} % modified
\newcommand{\fouhe}[2]{S_{#1}^{\{#2\}}} % Fourier-Hermite
% Cesaro-Laguerre kernel and means:
\newcommand{\claker}[5]{\mathcal{K}_{#1,#2}^{(#3)}%
  (#4,#5)} % Cesaro-Laguerre kernel
\newcommand{\clam}[3]{\sigma_{#1,#2}^{(#3)}} % Cesaro-Laguerre mean
% Hermite functions and polynomials:
\newcommand{\herfun}[2]{\mathcal{H}_{#1}^{\{#2\}}}
\newcommand{\herpol}[2]{H_{#1}^{\{#2\}}}
% Laguerre functions and polynomials:
\newcommand{\lagfun}[2]{\mathcal{L}_{#1}^{(#2)}}
\newcommand{\lagpol}[2]{L_{#1}^{(#2)}}
% Numeros combinatorios
\newcommand{\combi}[2]{A_{#1}^{#2}}
%
\newcommand{\I}{\infty}
\newcommand{\sgn}{\operatorname{sgn}}

\begin{document}
\title[Ces\`aro means of generalized Hermite expansions]
{Two-weight norm inequalities for the Ces\`aro means
of generalized Hermite expansions}

\author[\'O. Ciaurri]{\'Oscar Ciaurri}
\address{Departamento de Matem\'aticas y Computaci\'on,
Universidad de La Rioja, Edificio J.~L.~Vives,
Calle Luis de Ulloa s/n, 26004 Logro\~no, Spain}
\email{oscar.ciaurri@dmc.unirioja.es}

\author[J. L. Varona]{Juan~L.~Varona}
\address{Departamento de Matem\'aticas y Computaci\'on,
Universidad de La Rioja, Edificio J.~L.~Vives,
Calle Luis de Ulloa s/n, 26004 Logro\~no, Spain}
\email{jvarona@dmc.unirioja.es}
\urladdr{http://www.unirioja.es/dptos/dmc/jvarona/welcome.html}
\thanks{Research supported by grant BFM2003-06335-C03-03 of the DGI}

\keywords{Ces\`aro means, N{\o}rlund means, Hermite expansions, Hermite polynomials,
generalized Hermite expansions, generalized Hermite polynomials,
two-weight norm inequalities, weighted norm inequalities}
\subjclass[2000]{Primary 42C10}

\begin{abstract}
We prove two-weight norm inequalities for Ces\`aro means of generalized Hermite polynomial series
and for the supremum of these means. A result about weak boundedness and an almost everywhere
convergence result are also obtained.
\end{abstract}

%%%%%%%%%%%%%%%%%%%%%%%%%%%%%%%%%%%%%%%%%%%%%%%%%%%%%%%%%%%%%%%%%%%%%%%
\maketitle
%%%%%%%%%%%%%%%%%%%%%%%%%%%%%%%%%%%%%%%%%%%%%%%%%%%%%%%%%%%%%%%%%%%%%%%

%%%%%%%%%%%%%%%%%%%%%%%%%%%%%%%%%%%%%%%%%%%%%%%%%%%%%%%%%%%%%%%%%%%%%%%
\section{Introduction}
%%%%%%%%%%%%%%%%%%%%%%%%%%%%%%%%%%%%%%%%%%%%%%%%%%%%%%%%%%%%%%%%%%%%%%%

Given a real number $\mu > -1/2$, let $\{\herpol{n}{\mu}\}_{n=0}^\I$ be the
so-called generalized Hermite polynomials of order~$\mu$;
see~\cite[Ch.~V, \S\,2 (G), p.~156]{Ch}.
These polynomials are orthogonal on $L^2(\mathbb{R},|x|^{2\mu} e^{-x^2}\,dx)$,
and they are uniquely defined by requiring that the leading coefficient
of the polynomial $\herpol{n}{\mu}(x)$ be~$2^n$. 
The sequence of generalized Hermite functions
$\{\herfun{n}{\mu}\}_{n = 0}^{\I}$ is defined by
\[
  \herfun{n}{\mu}(x) :=
  2^{-n} \left(\Bigl[\frac{n}{2}\Bigr]!\right)^{-1/2}
    \Gamma\left(\Bigl[\frac{n+1}{2}\Bigr]
      +\mu+\frac{1}{2}\right)^{-1/2}
  |x|^{\mu} e^{-x^2/2} \herpol{n}{\mu}(x),
\]
where, as usual, $[\,\cdot\,]$ denotes the greatest
integer function.
The system $\{\herfun{n}{\mu}\}_{n=0}^\I$ is orthonormal on
$L^2(\mathbb{R},dx)$.
Of course, these polynomials and functions are generalizations
of the ``ordinary'' Hermite polynomials and functions, which correspond
to the case~$\mu=0$. They play an important role as
eigenfunctions of the Dunkl transform.

When we have an orthogonal system and a function, we can consider its Fourier series
with respect to the system. An interesting problem is the study of the convergence
of such Fourier expansions. If the system is complete, this is always true for
functions in the appropriate $L^2$ space, but not always for functions
in $L^p$, $p\ne2$; also, we cannot ensure almost everywhere convergence.
It is well known the relation between the uniform boundedness of partial sum operators and the
convergence of the series, a crucial fact in this problem. When studying the uniform
boundedness, extra weights can be added, and this leads to convergence in different weighted $L^p$
spaces. If the convergence of the Fourier series fails, another summation methods can be considered;
in particular, the convergence of Ces\`aro means.

These kind of questions have been widely studied for the classical
Hermite system~\cite{AW,M2,Po,MW2}, but not for the generalized Hermite system.
This is the aim of this paper. In particular, we are going to extend
the results of~\cite{MW2} concerning Ces\`aro means of Hermite expansions
to generalized Hermite expansions.

For a function $f$, let $\chem{\delta}{n}{\mu}(f,x)$ be the $n$th
Ces\`aro mean of order $\delta>0$ of the expansion of $f$ in orthonormalized
generalized Hermite functions $\{\herpol{n}{\mu}\}_{n=0}^\I$ of order $\mu>-1/2$.
Also, let us use $\|\cdot\|_p$ to denote the unweighted $L^p$ norm on $(-\I,\I)$.
We prove inequalities of the form
\begin{equation*}
  \sup_{n\ge0} \, \Bigl\| |x|^a(1+|x|)^{b-a}\chem{\delta}{n}{\mu}(f,x) \Bigr\|_p
  \le C \Bigl\| |x|^A(1+|x|)^{B-A}f(x) \Bigr\|_p
\end{equation*}
for $1 \le p \le \infty$ (Theorem~\ref{MCGWT1}), and
\begin{equation*}
  \Bigl\| |x|^a(1+|x|)^{b-a} \sup_{n\ge0} |\chem{\delta}{n}{\mu}(f,x)| \Bigr\|_p
  \le C \Bigl\||x|^A(1+|x|)^{B-A}f(x) \Bigr\|_p
\end{equation*}
for $1 < p \le \infty$ (Theorem~\ref{MCGW2}), where $C$ is independent of~$f$.
A result about weak boundedness for the supremum when $p=1$ (Theorem~\ref{WK}),
and an almost everywhere convergence result (Theorem~\ref{CAE}) are also proved.

An important point in the study of this question is to obtain
estimates for the kernel of $\chem{\delta}{n}{\mu}$.
This can be done in a direct way, such as it is done in~\cite{M}
for the case $\mu=0$ and $\delta=0$ (i.e., the Fourier series itself).
On the other hand, the kernel can be written in terms of Ces\`aro-Laguerre
kernels, and so previously known results for Laguerre can be applied.

If the order of summation, $\delta$, is an integer,
this is simple since the generalized Ces\`aro-Hermite kernel is a
linear combination of a fixed finite number of Ces\`aro-Laguerre
kernels. This fact was used for the first time in~\cite{Po},
for the case $\mu=0$ and $\delta=1$.
In \cite[Chapter~6]{V}, this is done for the general case $\mu>-1/2$,
both for $\delta=0$ and $\delta=1$.

For arbitrary $\delta>0$, obtaining the estimate for the kernel of
$\chem{\delta}{n}{\mu}$ is more complicated,
because the expression that relates Hermite kernels with Laguerre kernels
contains $n+1$ terms. When $\mu=0$, this decomposition is shown
in~\cite[Lemma~(3.8)]{MW2}; and the generalization for the arbitrary
$\mu > -1/2$ is not difficult.
Combining the decomposition with the precise estimates
for Laguerre kernels of~\cite{W}, an estimate for classical Hermite kernels
is given in~\cite[Theorem~(4.5)]{MW2}. Then, because of the similarity of the
estimate to the Laguerre case, the norm inequalities are established
using the methods and the results of~\cite{MW}.

In this paper, we extend the aforementioned studies to the general case $\mu>-1/2$.
But this is not the only purpose of the paper, but also to show a different approach:
we make use of the results in~\cite{MW}, but in a way different to the way in~\cite{MW2}.

Instead of using Ces\`aro means of generalized Hermite series, we will use a different
summation method, whose means we will denote by $\chemmod{\delta}{n}{\mu}$
(see~\eqref{stilde} for details). Actually, it will be a N{\o}rlund method,
according to the name given in~\cite[Chapter~IV]{H} (also Ces\`aro means are particular
cases of N{\o}rlund methods).
We will see that $\chemmod{\delta}{n}{\mu}$ can be decomposed as a sum of two
Ces\`aro-Laguerre means (see the formula~\eqref{twosum}); i.e, we have only two summands
in the relation between Laguerre and Hermite, instead of the $n+1$ summands that
appear when using Ces\`aro-Hermite. As a consequence of the fixed number of summands,
the uniform boundedness of $\chemmod{\delta}{n}{\mu}$ follows immediately from the
uniform boundedness of
Ces\`aro-Laguerre: contrary to~\cite{MW2}, it is not necessary the cumbersome process
of finding bounds for the kernel of~$\chemmod{\delta}{n}{\mu}$. Also, we completely
eliminate the part corresponding to the use of the bounds of the kernel to study
the uniform boundedness of the operators~$\chem{\delta}{n}{\mu}$ (this would require
to mimic the process in~\cite{MW}, as explained in~\cite{MW2}).

Finally, we prove that the uniform boundedness for~$\chemmod{\delta}{n}{\mu}$ imply
the uniform boundedness for the Ces\`aro-Hermite means~$\chem{\delta}{n}{\mu}$;
this is the most technical part of the paper.
In practice, due to the similarity to the Laguerre case, the norm inequalities,
necessity results and convergence results in this paper are essentially corollaries of
the results in~\cite{MW} for Laguerre expansions with parameter $\alpha=\mu-1/2$.
This happens because the hypotheses over the other parameter involved,
$\mu+1/2$, are weaker.

Throughout this paper $C$ will be a positive constant independent of $f$, $n$, $x$ and
$y$, but it assumes different values in different occurrences.

%%%%%%%%%%%%%%%%%%%%%%%%%%%%%%%%%%%%%%%%%%%%%%%%%%%%%%%%%%%%%%%%%%%%%%%
\section{Generalized Hermite in terms of Laguerre}
%%%%%%%%%%%%%%%%%%%%%%%%%%%%%%%%%%%%%%%%%%%%%%%%%%%%%%%%%%%%%%%%%%%%%%%

In the introduction, we have already described generalized
Hermite polynomials and functions. Let us now describe
Laguerre polynomials and functions. To clearly differentiate
between Hermite and Laguerre (in polynomials, functions, series,
Ces\`aro means, kernels, bounds, \dots), we will always use superscripts
$^{\{\cdot\}}$ to indicate the parameter for Hermite,
and superscripts $^{(\cdot)}$ for Laguerre parameters.

Given a real number $\alpha > -1$, let $\{\lagpol{n}{\alpha}\}_{n=0}^\I$
be the Laguerre polynomials of order~$\alpha$;
see, for instance,~\cite[Ch.~V, \S\,2 (A), p.~144]{Ch}.
These polynomials are orthogonal on $L^2((0,\I),x^\alpha e^{-x}\,dx)$,
and they are uniquely defined by requiring that the leading coefficient
of $\lagpol{n}{\alpha}(x)$ be~$(-1)^n/n!$.
The sequence of Laguerre functions
$\{\lagfun{n}{\alpha}\}_{n = 0}^{\I}$ is defined by
\[
  \lagfun{n}{\alpha}(x) :=
  {\Gamma(\alpha + 1)}^{-1/2}
    (\combi{n}{\alpha})^{-1/2} x^{\alpha/2}
    e^{-x/2} \lagpol{n}{\alpha}(x),
\]
where $\combi{n}{\alpha} = \binom{n+\alpha}{n}$. The system $\{\lagfun{n}{\mu}\}_{n=0}^\I$ is
orthonormal on $L^2((0,\I),dx)$.

It is interesting to note that, for Hermite polynomials,
only the case $\mu=0$ is ``classical'', according to the characterizations
of the classical orthogonal polynomials; see~\cite[Ch.~V, \S\,2 (D), p.~150]{Ch}.
However, Laguerre polynomials are considered as classical for every $\alpha>-1$,
although only the case $\alpha=0$ was originally studied by Laguerre.

The generalized Hermite polynomials are related to the Laguerre polynomials by the identities
\begin{align*}
\herpol{2n}{\mu}(x) &= (-1)^n 2^{2n}n!\,\lagpol{n}{\mu-1/2}(x^2),\\
\herpol{2n+1}{\mu}(x) &= (-1)^n 2^{2n+1}n!\,x\lagpol{n}{\mu+1/2}(x^2)
\end{align*}
(we must comment that in~\cite[Ch.~V, \S\,2 (G), p.~156]{Ch} there is a misprint in the second identity, the factor $x$ was omitted).
Then, each $\herfun{n}{\mu}$ can be expressed in terms of some
$\lagfun{n}{\alpha}$; namely, given a nonnegative integer $n$ and a real number $x$,
\begin{equation}
\label{h-lid}
\begin{aligned}
  \herfun{2n}{\mu}(x) &= (-1)^{n}|x|^{1/2}
    \lagfun{n}{\mu-1/2}(x^{2}),\\
  \herfun{2n + 1}{\mu}(x) &= (-1)^{n}\sgn(x)|x|^{1/2}
    \lagfun{n}{\mu+1/2}(x^{2}).
\end{aligned}
\end{equation}

The generalized Hermite expansion of a function $f$ is
\[
  \sum_{k=0}^\infty\herfun{k}{\mu}(x)
  \left(\int_{-\I}^\I f(y)\herfun{k}{\mu}(y)\,dy\right)
\]
provided that the integrals (i.e., the Fourier coefficients) exist.
For $\delta>0$, the $n$th $(C,\delta)$-Ces\`aro
mean of this expansion is
\[
  \chem{\delta}{n}{\mu}(f,x)
  := \frac{1}{\combi{n}{\delta}} \sum_{k=0}^n\combi{n-k}{\delta}
  \herfun{k}{\mu}(x) \left(\int_{-\I}^\I f(y)\herfun{k}{\mu}(y)\,dy\right).
\]
It follows that
\begin{equation*}
  \chem{\delta}{n}{\mu}(f,x) = \int_{-\I}^{\I}f(y)
    \cheker{\delta}{n}{\mu}{x}{y}\,dy,
\end{equation*}
where
\begin{equation*}
  \cheker{\delta}{n}{\mu}{x}{y}
  := \frac{1}{\combi{n}{\delta}} \sum_{k = 0}^{n}
  \combi{n-k}{\delta} \herfun{k}{\mu}(x)\herfun{k}{\mu}(y).
\end{equation*}
Similarly, the $n$th ($C,\delta$)-Ces\`aro mean for a Laguerre expansion satisfies
\begin{equation*}
  \clam{\delta}{n}{\alpha}(f,x)
  = \int_{0}^{\I}f(y) \claker{\delta}{n}{\alpha}{x}{y}\,dy,
\end{equation*}
where
\begin{equation*}
  \claker{\delta}{n}{\alpha}{x}{y}
  := \frac{1}{\combi{n}{\delta}} \sum_{k = 0}^{n}
  \combi{n-k}{\delta} \lagfun{k}{\alpha}(x)\lagfun{k}{\alpha}(y).
\end{equation*}

The main tool to prove the results in this paper will be the use of another summation method.
For a generalized Hermite expansion, the $n$th $(\widetilde{C},\delta)$ mean is defined by
\begin{equation}
\label{stilde}
  \chemmod{\delta}{n}{\mu}(f,x)
  := \frac{1}{\combi{[\frac{n}2]}{\delta}}
  \sum_{k=0}^n\combi{[\frac{n-k}2]}{\delta}
  \herfun{k}{\mu}(x) \left(\int_{-\I}^\I f(y)\herfun{k}{\mu}(y)\,dy\right).
\end{equation}
The point of the $(\widetilde{C},\delta)$ summation method is that we can express it easily
in terms of the Ces\`{a}ro means of two Laguerre expansions. So, norm inequalities for
$\chemmod{\delta}{n}{\mu}$ can be deduced from the results in \cite{MW} for
Ces\`{a}ro-Laguerre expansions. This will be done in Proposition~\ref{Propo1}. Moreover,
as we will show in Proposition~\ref{Propo2}, norm inequalities for $\chem{\delta}{n}{\mu}$
follow from the norm inequalities for~$\chemmod{\delta}{n}{\mu}$.

\begin{Prop}
\label{Propo1}
Let $1\le p \le \I$, $\mu >-1/2$ and $\delta>0$.
\begin{enumerate}
\item[a)] If
\[
  \sup_{n\ge 0} \left\|w(x)\int_{-\I}^{\I}f(y)|xy|^{1/2}
    \claker{\delta}{n}{\mu\pm 1/2}{x^2}{y^2}\,dy\right\|_{p}
  \le C \|W(x)f(x)\|_{p}
\]
then
\[
  \sup_{n\ge 0} \left\|w(x)\chemmod{\delta}{n}{\mu}(f,x)\right\|_p
  \le C \|W(x)f(x)\|_p.
\]
\item[b)] If
\[
  \left\|w(x)\sup_{n\ge 0} \left|\int_{-\I}^{\I}f(y)|xy|^{1/2}
    \claker{\delta}{n}{\mu\pm 1/2}{x^2}{y^2}\,dy\right|\right\|_{p}
  \le C \|W(x)f(x)\|_{p}
\]
then
\[
  \left\|w(x)\sup_{n\ge 0} \left|\chemmod{\delta}{n}{\mu}(f,x)\right|\right\|_p
  \le C \|W(x)f(x)\|_p.\]
\end{enumerate}
\end{Prop}
The proof of this result follows immediately from the identity
\begin{multline}
\label{twosum}
  \chemmod{\delta}{n}{\mu}(f,x)
  = \int_{-\I}^{\I}f(y)|xy|^{1/2}
    \claker{\delta}{[\frac{n}2]}{\mu-1/2}{x^2}{y^2}\,dy \\
  + \int_{-\I}^{\I}f(y)\sgn(xy)|xy|^{1/2}
    \claker{\delta}{[\frac{n-1}2]}{\mu+1/2}{x^2}{y^2}\,dy,
\end{multline}
which can be easily obtained by using~\eqref{h-lid}.

\begin{Rem}
\label{equiv}
Actually, we can prove the reverse condition both in parts (a) and~(b).
For this, it suffices to take even and odd functions.
Then, one of the two summands in~\eqref{twosum} vanishes,
and so the boundedness of $\chemmod{\delta}{n}{\mu}$ is equivalent
to the boundedness of Ces\`aro-Laguerre means.
In particular, this would ensure that, in the situation of Theorem~\ref{MCGWT1},
the sufficient conditions are also necessary for the uniform boundedness
of~$\chemmod{\delta}{n}{\mu}$,
because this is what happens in the Laguerre case~\cite{MW}.
\end{Rem}

\begin{Prop}
\label{Propo2} Let $1\le p \le \I$, $\mu>-1/2$ and $\delta>0$.
\begin{enumerate}
\item[a)] If
\[
  \sup_{n\ge 0} \left\|w(x)\chemmod{\delta}{n}{\mu}(f,x)\right\|_{p}
  \le C \|W(x)f(x)\|_{p}
\]
then
\[
  \sup_{n\ge 0} \left\|w(x)\chem{\delta}{n}{\mu}(f,x)\right\|_p
  \le C \|W(x)f(x)\|_p.
\]
\item[b)] If
\[
  \left\|w(x)\sup_{n\ge 0} \left|\chemmod{\delta}{n}{\mu}(f,x)\right|\right\|_{p}
  \le C \|W(x)f(x)\|_{p}
\]
then
\[
  \left\|w(x)\sup_{n\ge 0} \left|\chem{\delta}{n}{\mu}(f,x)\right|\right\|_p
  \le C \|W(x)f(x)\|_p.
\]
\end{enumerate}
\end{Prop}

\begin{proof}
The proof of these results relies on the ideas of \cite[\S\,4.3]{H} about
the \emph{inclusion} of different N{\o}rlund methods. Let us begin
writing $\chem{\delta}{n}{\mu}$ in terms of $\chemmod{\delta}{n}{\mu}$.
By using the identity
$\combi{m}{\delta}-\combi{m-1}{\delta} = \combi{m}{\delta-1}$,
it is clear that
\[
  \chem{\delta}{n}{\mu}(f,x)
  = \frac{1}{\combi{n}{\delta}} \sum_{k=0}^{n}\combi{n-k}{\delta-1} \fouhe{k}{\mu}(f,x),
\]
where
\[
  \fouhe{k}{\mu}(f,x) = \sum_{i=0}^k \herfun{i}{\mu}(x)
   \left(\int_{-\I}^\I f(y)\herfun{i}{\mu}(y)\,dy\right)
\]
are the partial sums of the Fourier series.
In a similar way, we can show that
\[
  \chemmod{\delta}{n}{\mu}(f,x) =
  \begin{cases} \displaystyle
    \frac{1}{\combi{m}{\delta}} \sum_{k=0}^m \combi{m-k}{\delta-1} \fouhe{2k}{\mu}(f,x),
      &\text{if $n=2m$,}\\[4mm]
    \displaystyle
    \frac{1}{\combi{m}{\delta}}\sum_{k=0}^m \combi{m-k}{\delta-1} \fouhe{2k+1}{\mu}(f,x),
      &\text{if $n=2m+1$.}
  \end{cases}
\]
From this point on, we will consider $n=2m$. The case $n=2m+1$ is similar. It is easy to verify that
\begin{align*}
  \chem{\delta}{2m}{\mu}(f,x)
  &= \frac{1}{\combi{2m}{\delta}} \sum_{j=0}^{m}\combi{2(m-j)}{\delta-1} \fouhe{2j}{\mu}(f,x)
  + \frac{1}{\combi{2m}{\delta}} \sum_{j=0}^{m-1}\combi{2(m-j)-1}{\delta-1} \fouhe{2j+1}{\mu}(f,x) \\
  &=: \chem{\delta}{2m}{\mu,1}(f,x) + \chem{\delta}{2m}{\mu,2}(f,x).
\end{align*}
Now, we claim that
\begin{align}
\label{ec:Id1}
  \chem{\delta}{2m}{\mu,1}(f,x)
  &= \frac{1}{\combi{2m}{\delta}}\sum_{j=0}^m
    \combi{j}{\delta} \binom{\delta}{2(m-j)} \chemmod{\delta}{2j}{\mu}(f,x)\\
\intertext{and}
\label{ec:Id2}
  \chem{\delta}{2m}{\mu,2}(f,x)
  &= \frac{1}{\combi{2m}{\delta}} \sum_{j=0}^{m-1}
   \combi{j}{\delta}\binom{\delta}{2(m-j)-1}\chemmod{\delta}{2j+1}{\mu}(f,x).
\end{align}
Then, the proof can be concluded using two facts:
for a constant $C$ independent of $m$ and~$0 \le j \le m$,
the binomial coefficients satisfy $\combi{j}{\delta}\le C \combi{2m}{\delta}$;
and, finally, $\sum_{j=0}^\I |\binom{\delta}{j}| < \I$.

Now, let us check our claim. We only prove~\eqref{ec:Id1},
because the identity~\eqref{ec:Id2} follows in the same way.
Taking the power series (which are absolutelly convergent for $|t|$ small enough)
\[
  s(t) = \sum_{m=0}^\I \fouhe{2m}{\mu}(f,x)t^m,\quad p(t)
  = \sum_{m=0}^\I \combi{2m}{\delta-1}t^m\quad
  \text{and}\quad
  q(t) = \sum_{m=0}^\I \combi{m}{\delta-1}t^m,
\]
it is clear that
\[
  s(t)p(t) = \sum_{m=0}^\I \combi{2m}{\delta} \chem{\delta}{2m}{\mu,1}(f,x)t^m
  \quad\text{and}\quad
  s(t)q(t) = \sum_{m=0}^\I \combi{m}{\delta} \chemmod{\delta}{2m}{\mu}(f,x)t^m.
\]
Moreover,
\[
  p(t) = \frac{(1+\sqrt{t})^{\delta}+(1-\sqrt{t})^{\delta}}{2(1-t)^\delta},
  \quad q(t) = \frac{1}{(1-t)^\delta}
\]
and
\[
  \frac{p(t)}{q(t)} = \frac12((1+\sqrt{t})^\delta + (1-\sqrt{t})^\delta)
  = \sum_{m=0}^\I \binom{\delta}{2m}t^m.
\]
In this way,
\begin{align*}
  \sum_{m=0}^\I \combi{2m}{\delta} \chem{\delta}{2m}{\mu,1}(f,x)t^m
  &= \frac{p(t)}{q(t)}s(t)q(t)
  = \sum_{m=0}^\I \binom{\delta}{2m}t^m
    \cdot \sum_{m=0}^\I \combi{\delta}{m} \chemmod{\delta}{2m}{\mu}(f,x)t^m \\
  &= \sum_{m=0}^\I \left(\sum_{j=0}^m \combi{j}{\delta}\binom{\delta}{2(m-j)}
    \chemmod{\delta}{2j}{\mu}(f,x)\right)t^m
\end{align*}
and so we have~\eqref{ec:Id1}.
\end{proof}

\begin{Rem}
\label{sst}
In general, the converse results are not true. If we try to prove them,
a problem arises: the series $\sum_{j=0}^\I |\binom{-\delta}{j}|$ does not converge.
\end{Rem}


%%%%%%%%%%%%%%%%%%%%%%%%%%%%%%%%%%%%%%%%%%%%%%%%%%%%%%%%%%%%%%%%%%%%%%%
\section{Norm inequalities and convergence results}
%%%%%%%%%%%%%%%%%%%%%%%%%%%%%%%%%%%%%%%%%%%%%%%%%%%%%%%%%%%%%%%%%%%%%%%

Here, we establish the main results of the paper.
Most of their proofs are immediate by applying~\cite{MW} together
with Propositions~\ref{Propo1} and~\ref{Propo2}; only the necessary conditions will require some extra comments.
Analogs of other results in~\cite{MW} or~\cite{MW2} could be obtained similarly;
in particular, the theorems corresponding to the case $a=A=b=B=r$.

\begin{Def}
\label{DNp} Let $1\le p\le\infty$, $\mu>-1/2$ and $\delta>0$. We say that parameters
$(a,b,A,B,\mu,\delta)$ satisfy the $H\!N_p$ conditions if
\begin{align}
\label{Na1} a &\ge -\delta-1/p,\\
\label{Na2} a &> -\mu-1/p && (\ge \textup{if } p=\infty),\\
\label{NAa} A-a &\le 0,\\
\label{NA1} A &\le 1+\delta-1/p,\\
\label{NA2} A &< 1+\mu-1/p && (\le \textup{if } p=1),\\
\label{NaB1} a+B &\ge -1-2\delta-2/(3p),\\
\label{NaB2} a+B &\ge -2\delta-2/p,\\
\label{NAb1} A+b &\le 2+2\delta-2/p,\\
\label{NAb2} A+b &\le 5/3+2\delta-2/(3p),\\
\label{Nb1} b &\le 1+2\delta-1/p,\\
\label{Nb2} b &\le 2/3+2\delta+1/(3p),\\
\label{NbB1} b-B &\le 1+2\delta-4/(3p),\\
\label{NbB2} b-B &\le 0,\\
\label{NbB3} b-B &\le -1/3+2\delta+4/(3p),\\
\label{NB1} B &\ge -1-2\delta+1/(3p),\\
\label{NB2} B &\ge -2\delta-1/p,
\end{align}
and in at least one of each of the following pairs
the inequality is strict:
\eqref{Na1} and \eqref{NAa} except for $p=1$,
\eqref{Na1} and \eqref{NaB2},
\eqref{NAa} and \eqref{NA1} except for $p=\I$,
\eqref{NA1} and \eqref{NAb1},
\eqref{NaB1} and \eqref{NaB2},
\eqref{NaB1} and \eqref{NB1},
\eqref{NaB2} and \eqref{NB2} except for $p=\I$,
\eqref{NAb1} and \eqref{NAb2},
\eqref{NAb1} and \eqref{Nb1} except for $p=1$,
\eqref{NAb2} and \eqref{Nb2},
\eqref{Nb1} and \eqref{Nb2},
\eqref{Nb2} and \eqref{NbB3},
\eqref{NbB1} and \eqref{NB1},
\eqref{NB1} and \eqref{NB2}.
\end{Def}

\begin{Theo}
\label{MCGWT1}
Let $1\le p\le\infty$, $\mu>-1/2$, $\delta>0$, and suppose $(a,b,A,B,\mu,\delta)$
satisfy the $H\!N_{p}$ conditions. Then
\begin{equation}
\label{MCGW1}
  \sup_{n\ge0} \left\| |x|^a(1+|x|)^{b-a}\chem{\delta}{n}{\mu}(f,x) \right\|_p
  \le C \left\| |x|^A(1+|x|)^{B-A}f(x) \right\|_p
\end{equation}
with $C$ independent of~$f$.
Conversely, let us suppose that~\eqref{MCGW1} holds; thus, if $\mu\le \delta$, then
$(a,b,A,B,\mu,\delta)$ satisfy the $H\!N_p$ conditions; and, if $\mu>\delta$, then
$(a,b,A,B,\mu,\delta)$ satisfy the $H\!N_p$ conditions except for, perhaps,~\eqref{Na1},
\eqref{NA1} and their pair conditions.
\end{Theo}

For $p\ne\infty$, it is a corollary that,
under the conditions of the previous theorem, we have
$\displaystyle{\lim_{n\to\infty}\chem{\delta}{n}{\mu}f=f}$
in the $L^p((-\I,\I),|x|^{ap}(1+|x|)^{(b-a)p}\,dx)$-norm for every
$f \in L^p((-\I,\I),|x|^{Ap}(1+|x|)^{(B-A)p}\,dx)$.

\begin{Def}
\label{SNp}
Let $1\le p\le\infty$, $\mu>-1/2$ and $\delta>0$. We say that parameters
$(a,b,A,B,\mu,\delta)$ satisfy the $H\!S_p$ conditions if they satisfy inequalities
\eqref{Na2}--\eqref{NaB2}, \eqref{NbB2}, \eqref{NB1}--\eqref{NB2},
\begin{align}
\label{Sa} a &> -\delta-1/p && (\ge\textup{if } p=\I),\\
\label{SAb} A+b &\le 5/3+2\delta-2/p,\\
\label{Sb} b &< 2/3+2\delta-1/p && (\le\textup{if } p=\I),\\
\label{SbB} b-B &\le -1/3 + 2\delta,
\end{align}
and in at least one of each of the following pairs
the inequality is strict:
\eqref{NAa} and \eqref{NA1} except for $p=\I$,
\eqref{NAa} and \eqref{NA2},
\eqref{NAa} and \eqref{Sa},
\eqref{NA2} and \eqref{SAb},
\eqref{NaB1} and \eqref{NaB2},
\eqref{NaB1} and \eqref{NB1},
\eqref{NaB2} and \eqref{NB2} except for $p=\I$,
\eqref{NaB2} and \eqref{Sa},
\eqref{NB1} and \eqref{NB2},
\eqref{NB1} and \eqref{SbB} for $p=1$,
\eqref{SAb} and \eqref{Sb},
\eqref{Sb} and \eqref{SbB}.
\end{Def}

\begin{Theo}
\label{MCGW2}
Let $1<p\le\infty$, $\mu>-1/2$, $\delta>0$, and suppose $(a,b,A,B,\mu,\delta)$ satisfy
the $H\!S_p$ conditions. Then
\begin{equation}
\label{SCGW1}
  \left\| |x|^a(1+|x|)^{b-a}\sup_{n\ge0}|\chem{\delta}{n}{\mu}(f,x)| \right\|_p
  \le C \left\| |x|^A(1+|x|)^{B-A}f(x) \right\|_p
\end{equation}
with $C$ independent of~$f$.
\end{Theo}

Note that since \eqref{SCGW1} implies \eqref{MCGW1}, the necessary conditions of
Theorem~\ref{MCGWT1} are also necessary for \eqref{SCGW1}. On the other hand, the condition
\eqref{SbB} is not necessary: following the method of~\cite[\S\,10]{MW}, some hypotheses that
guarantee \eqref{SCGW1} with $b-B>2\delta-1/3$ can be found.

\begin{Theo}
\label{WK} If $\mu>-1/2$, $\delta>0$, $(a,b,A,B,\mu,\delta)$ satisfy the $H\!S_1$ conditions and
$E_{\lambda}$ is the set where $|x|^a(1+|x|)^{b-a} \sup_{n\ge0}
\bigl(\bigl|\chem{\delta}{n}{\mu}(f,x)\bigr|\bigr) > \lambda$, then
\[
  \left|E_\lambda\right|
  \le (C/\lambda) \left\| |x|^A(1+|x|)^{B-A}f(x) \right\|_1
\]
holds with $C$ independent of~$f$ and~$\lambda$.
\end{Theo}

\begin{Theo}
\label{CAE} If $1\le p\le\infty$, $\mu>-1/2$, $\delta>0$, \eqref{NA1}, \eqref{NA2}, \eqref{NB1}
and \eqref{NB2} are satisfied with equality in at most one of \eqref{NB1} and \eqref{NB2}, and
\[
  \left\| |x|^A(1+|x|)^{B-A}f(x)\right\|_p < \infty,
\]
then
$\displaystyle{\lim_{n\to\infty}\chem{\delta}{n}{\mu}(f,x)=f(x)}$
for almost every $x \in \mathbb{R}$.
\end{Theo}

Theorem~\ref{CAE} is proved by choosing an $a$ large enough and a $b$ small enough that $a$, $A$,
$b$ and $B$ satisfy the conditions of Theorem~\ref{MCGW2} if $p>1$ or Theorem~\ref{WK} if $p=1$.
The conclusions of those theorems then imply the almost everywhere convergence by a standard
argument.

To prove the sufficiency of the conditions in Theorem~\ref{MCGWT1}, applying
Propositions~\ref{Propo1} and~\ref{Propo2}, and using that the kernel and the weight
functions are even in both $x$ and $y$, it is enough to show that
\begin{multline*}
  \int_0^\I x^{ap}(1+x)^{(b-a)p} \left|\int_0^\I f(y)(xy)^{1/2}
    \claker{\delta}{n}{\mu\pm 1/2}{x^2}{y^2}\,dy\right|^p\,dx \\
  \le C \int_0^\I x^{Ap}(1+x)^{(B-A)p}|f(x)|^p\,dx.
\end{multline*}
With the change of variables $x=\sqrt{z}$, $y=\sqrt{u}$, and taking
$g(u) = u^{-1/4}|f(\sqrt{u})|$, this is equivalent to
\begin{multline*}
  \int_0^{\I}z^{(2a+1)p/4-1/2}(1+z)^{(b-a)p/2}|
    \clam{\delta}{n}{\mu\pm 1/2}(g,z)|^p\,dz \\
  \le C \int_0^\I z^{(2A+1)p-1/2}(1+z)^{(B-A)p/2}|g(z)|^p\,dz.
\end{multline*}
But this inequality holds provided that the parameters
\begin{equation}
\label{prm}
  \left( \frac{a}{2}+\frac{1}{4}-\frac{1}{2p},
  \ \frac{b}{2}+\frac{1}{4}-\frac{1}{2p},
  \ \frac{A}{2}+\frac{1}{4}-\frac{1}{2p},
  \ \frac{B}{2}+\frac{1}{4}-\frac{1}{2p},
  \ \mu\pm\frac{1}{2},\ \delta \right)
\end{equation}
satisfy the corresponding $N_p$ conditions for Laguerre, defined on~\cite[p.~1125--1126]{MW}. It
is now a routine procedure to show that these parameters satisfy the $N_p$ conditions for Laguerre
if $(a,b,A,B,\mu,\delta)$ satisfy the hypotheses of Theorem~\ref{MCGWT1}, i.e., our $H\!N_p$
conditions of Definition~\ref{DNp}. This completes the sufficiency part of Theorem~\ref{MCGWT1}.

Theorems~\ref{MCGW2} and~\ref{WK} are proved in the same way as the sufficiency proof of
Theorem~\ref{MCGWT1}. The same change of variables will reduce the proof to the sufficiency part
of Theorems~(2.30) and~(2.31) of~\cite{MW}, and it is simple to show that the
resulting parameters \eqref{prm} satisfy the corresponding $S_p$ conditions for Laguerre (defined
on~\cite[p.~1126]{MW}).

Now, let us analyze the necessity of the conditions in Theorem~\ref{MCGWT1}.
First, it is clear that the orthogonal functions must satisfy
$\herfun{n}{\mu}(x)|x|^a(1+|x|)^{b-a} \in L^p(\mathbb{R},dx)$;
this is equivalent to~\eqref{Na2}.
Second, we must ensure that the Fourier coefficients exist for every
function $f$ such that $f(x)|x|^A(1+|x|)^{B-A} \in L^p(\mathbb{R},dx)$;
by duality, this is equivalent to
$\herfun{n}{\mu}(x)|x|^{-A}(1+|x|)^{A-B} \in L^{p'}(\mathbb{R},dx)$
(being $1/p+1/p'=1$), and so we get the necessity of~\eqref{NA2}.

In addition, for a fixed $\delta>0$ and $r>0$, Theorem~\ref{CAE} implies that
$|\chem{\delta}{n}{\mu}(\chi_{[r,2r]},x)|$ converges almost everywhere to
$\chi_{[r,2r]}(x)$.  From Fatou's lemma and \eqref{MCGW1} it follows
that
\[
  \left\| x^a(1+x)^{b-a}\chi_{[r,2r]}(x) \right\|_p
  \le C \left\| x^A(1+x)^{B-A}\chi_{[r,2r]}(x) \right\|_p;
\]
\eqref{NAa} and \eqref{NbB2} follow from this.  Next, a standard argument as
given on~\cite[p.~1141]{MW} or~\cite[p.~113]{Th} shows that \eqref{MCGW1}
implies
\[
  \left\| |x|^a(1+|x|)^{b-a}\herfun{n}{\mu}(x) \right\|_p
  \left\| |x|^{-A}(1+|x|)^{A-B}\herfun{n}{\mu}(x) \right\|_{p'}
  \le C(n+1)^{\delta}.
\]
The necessity of the rest of the $H\!N_p$ conditions in Theorem~\ref{MCGWT1}, except~\eqref{Na1},
\eqref{NA1} and the pair restrictions for these inequalities, follow from this and the following
lemma:

\begin{Lem}
\label{SL}
Let $1\le p\le\infty$, $\mu>-1/2$ and $n\ge2$. Then,
\begin{multline*}
  \left\||x|^a(1+|x|)^{b-a}\herfun{2n}{\mu}(x)\right\|_p \\
  \ge C \left(n^{-1/4} + n^{-a/2-1/4-1/(2p)}
    + n^{b/2-1/4+1/(2p)} + n^{b/2-1/12-1/(6p)}\right).
\end{multline*}
Moreover, if $a=-1/p$ or $b=-1/p$, then
\[
  \left\||x|^a(1+|x|)^{b-a}\herfun{2n}{\mu}(x)\right\|_p
  \ge C n^{-1/4}(\log n)^{1/p};
\]
and, for $p=4$, we have
\[
  \left\||x|^a(1+|x|)^{b-a}\herfun{2n}{\mu}(x)\right\|_4
  \ge C n^{b/2-1/8}(\log n)^{1/4}.
\]
\end{Lem}

To prove Lemma~\ref{SL}, use
$\herfun{2n}{\mu}(x) = (-1)^n\sqrt{|x|}\lagfun{n}{\mu-1/2}(x^2)$ and
make a change of variables to show that
\[
  \left\||x|^a(1+|x|)^{b-a}\herfun{2n}{\mu}(x)\right\|_p
  \ge C\left\| x^{a/2+1/4-1/(2p)}(1+x)^{(b-a)/2}
    \lagfun{n}{\mu-1/2}(x) \right\|_p,
\]
where the norm on the right is taken over $[0,\infty)$. Then, the result
follows immediately from~\cite[Lemma~(7.2), p.~1142]{MW}.

To complete the proof of Theorem~\ref{MCGWT1}, it is enough to prove that~\eqref{Na1}
and~\eqref{NA1} are also necessary when $\mu\le\delta$. This is clear because they are
implied by~\eqref{Na2} and~\eqref{NA2}. In the case $\mu>\delta$, by applying
Lemma~\ref{SL}, instead of the conditions~\eqref{Na1} and~\eqref{NA1}, we get the
necessity of $a \ge -2\delta-1-1/p$ and $A \le 2\delta+2-1/p$, that are weaker.

Finally, let us comment the necessity (or not) of conditions \eqref{Na1} and~\eqref{NA1} (and the
pair restrictions involving them). When, in~\cite{MW2}, Ces\`aro-Laguerre series are studied, the
proof of the necessity of the conditions corresponding to \eqref{Na1} and~\eqref{NA1} is based on
the lower bounds for $\claker{\delta}{n}{\alpha}{x}{y}$ of \cite[Lemma~(8.1)]{MW} that, moreover,
are strongly dependent on~\cite{W}. As stated, it does not seem possible to apply them to
get lower bounds for $\cheker{\delta}{n}{\mu}{x}{y}$. However, it seems reasonable that these lower
bounds exist. It seems feasible to find them, but this would not be a direct consequence of the
Laguerre result; instead, this would require to reproduce a big part of~\cite{W,MW}, which is
outside of the purposes of the paper.

%%%%%%%%%%%%%%%%%%%%%%%%%%%%%%%%%%%%%%%%%%%%%%%%%%%%%%%%%%%%%%%%%%%%%%%

\begin{thebibliography}{99}

\bibitem{AW} R. Askey and S. Wainger,
{Mean convergence of expansions in Laguerre and Hermite series},
\textit{Amer. J. Math.}
\textbf{87} (1965), 695--708.

\bibitem{Ch} T. S. Chihara,
\textit{An Introduction to Orthogonal Polynomials},
Gordon and Breach, 1978.

\bibitem{H} G. H. Hardy,
\textit{Divergent Series},
Oxford at the Clarendon Press, Oxford, 1949.

\bibitem{M} B. Muckenhoupt,
{Asymptotic forms for Laguerre polynomials},
\textit{Proc. Amer. Math. Soc.}
\textbf{24} (1970), 288--292.

\bibitem{M2} B. Muckenhoupt,
{Mean convergence of Hermite and Laguerre series. I},
\textit{Trans. Amer. Math. Soc.}
\textbf{147} (1970), 419--431;
{II}, \textit{ibid.} \textbf{147} (1970), 433--460.

\bibitem{MW} B. Muckenhoupt and D. W. Webb,
{Two-weight norm inequalities for Ces\`aro means
of Laguerre expansions},
\textit{Trans. Amer. Math. Soc.}
\textbf{353} (2001), 1119--1149.

\bibitem{MW2} B. Muckenhoupt and D. W. Webb,
{Two-weight norm inequalities for the Ces\`aro means
of Hermite expansions},
\textit{Trans. Amer. Math. Soc.}
\textbf{354} (2002), 4525--4537.

\bibitem{Po} E. Poiani,
{Mean Ces\`aro summability of Laguerre and Hermite series},
\textit{Trans. Amer. Math. Soc.}
\textbf{173} (1972), 1--31.

\bibitem{Th} S. Thangavelu,
\textit{Lectures on Hermite and Laguerre Expansions},
Princeton Univ. Press, 1993.

\bibitem{V} J. L. Varona,
\textit{Convergencia en $L^p$ con pesos de la serie de Fourier
respecto de algunos sistemas ortogonales},
Doctoral Dissertation,
Publ. Sem. Mat. Garc\'{\i}a de Galdeano,
Ser.~II, no.~22, Zaragoza (Spain), 1989.

\bibitem{W} D. W. Webb,
{Pointwise estimates of the moduli of Ces\`aro-Laguerre kernels},
\textit{preprint}.

\end{thebibliography}

\end{document}
